% !TeX encoding = UTF-8
% Français de France. Charger un seul archétype depuis le document principal.
\newcommand{\cvsummary}{Ingénieur en mécanique et mécatronique formé à NYU Abu Dhabi, avec une expérience des systèmes mécaniques aux Émirats arabes unis et en Europe. Expérience de terrain en amélioration de l'efficacité énergétique industrielle, développement de prototypes éoliens et appui aux projets d'infrastructures hydrauliques.}


% --- COMPÉTENCES ---
\newcommand{\skills}{
    \section*{COMPÉTENCES TECHNIQUES}
    {\small
    \textbf{Énergie et développement durable :} Éolien flottant en mer (stabilité hydrostatique des plateformes semi-submersibles, validation en bassin à houle), aérodynamique des éoliennes à axe vertical (VAWT ; brevet en instance), systèmes thermiques (prototypage TEC et CVC), compensation du facteur de puissance et analyse de postes haute tension, analyse du marché du captage du carbone (CCUS), écoconception, métrologie de la mesure carbone.

    \vspace{1pt}
    \textbf{Simulation et analyse :} ANSYS (Mechanical, Fluent, AQWA), COMSOL Multiphysics, analyse par éléments finis (FEA), mécanique des fluides numérique (CFD), analyses thermiques et multiphysiques, modélisation dynamique sous MATLAB/Simulink.

    \vspace{1pt}
    \textbf{CAO et conception :} SolidWorks (CSWA), CATIA V5, Siemens NX, CREO, MoldFlow, Fusion 360, AutoCAD Electrical, conception paramétrique 3D, analyse structurelle et cinématique, cotation dimensionnelle et géométrique (GD\&T).

    \vspace{1pt}
    \textbf{Données et SIG :} PostgreSQL/PostGIS, Python, MATLAB, chaînes de traitement de données SQL, intégration de données géospatiales (SIG d'infrastructures hydrauliques), visualisation de données.

    \vspace{1pt}
    \textbf{Commande et instrumentation :} Commande PID, fusion de données inertielles (IMU), Arduino, automates programmables, capteurs et instrumentation, intégration matériel-logiciel.

    \vspace{1pt}
    \textbf{Gestion de projets :} MS Project, Jira, Notion, coordination pluridisciplinaire, rapports techniques clients, Lean Six Sigma White Belt.

    \vspace{1pt}
    \textbf{Langues :} Espagnol langue maternelle, anglais C2, français C2, italien A1, allemand A1.
    }
}



% --- PROJETS ---

\newcommand{\missiletitle}{Simulateur de dynamique d'interception de missiles}
\newcommand{\missilelocation}{Besançon, France}
\newcommand{\missiledescOne}{Développement sous MATLAB d'une simulation 3D d'interception de missiles avec loi de guidage par navigation proportionnelle, intégration numérique et visualisation des trajectoires en temps réel avec détection d'impact.}

\newcommand{\vawttitle}{VAWT -- Récupération d'énergie éolienne urbaine}
\newcommand{\vawtdates}{Automne 2024 - Printemps 2025}
\newcommand{\vawtlocation}{Abou Dabi, ÉAU}
\newcommand{\vawtdescOne}{Développement d'un récupérateur d'énergie éolienne urbaine faisant l'objet d'une demande de brevet déposée sur la captation du vent généré par le trafic, grâce à une éolienne à axe vertical dimensionnée pour les profils de rafales en bord de route.}
\newcommand{\vawtdescTwo}{Validation des performances avant fabrication par comparaison des résultats CFD d'ANSYS Fluent avec des données de vent mesurées, en simulant le comportement du rotor dans des régimes réalistes de vent généré par le trafic.}

\newcommand{\fsaetitle}{Formula SAE -- Suspension et châssis en fibre de carbone}
\newcommand{\fsaedates}{Automne 2024 -- Aujourd'hui}
\newcommand{\fsaelocation}{Abou Dabi, ÉAU}
\newcommand{\fsaedescOne}{Conception et simulation d'un ensemble suspension et châssis en fibre de carbone sous SolidWorks et ANSYS, avec performances structurelles vérifiées par éléments finis.}
\newcommand{\fsaedescTwo}{Réglage de la dynamique du véhicule par études cinématiques sous SolidWorks, avec analyse des courbes de carrossage, de roulis et de rapport de mouvement sur le débattement de la suspension.}

\newcommand{\tactitle}{Mini-réfrigérateur -- Système thermoélectrique TEC à deux étages}
\newcommand{\tacdates}{Printemps 2025}
\newcommand{\taclocation}{Abou Dabi, ÉAU}
\newcommand{\tacdescOne}{Suppression complète de la charge de fluide frigorigène grâce à une conception de refroidissement Peltier à deux étages, à l'état solide et sans fluide de travail, fondée sur des modules TEC.}
\newcommand{\tacdescTwo}{Gestion de la chaleur résiduelle par conception et impression 3D des dissipateurs et de l'ensemble de la structure de support, avec suivi de leurs performances thermiques.}

\newcommand{\orbitaltitle}{OrbitalGrid -- Système SIG pour applications spatiales}


\newcommand{\orbitaldates}{06/2024 - Aujourd'hui}
\newcommand{\orbitallocation}{Abou Dabi, ÉAU}
\newcommand{\orbitaldescOne}{Développement d'une base de données spatiale pour la cartographie aérospatiale, avec couches de numérisation 3D et de données satellitaires interrogeables sous PostgreSQL, grâce à la conception du schéma et des chaînes de chargement.}
\newcommand{\orbitaldescTwo}{Mise en place d'analyses croisées entre géométries numérisées et images satellitaires grâce à la normalisation des systèmes de coordonnées et de l'indexation spatiale.}

\newcommand{\softrobtitle}{Exosquelette robotique pour la rééducation après un AVC}
\newcommand{\softrobdates}{04/2024 - 05/2024}
\newcommand{\softroblocation}{Abou Dabi, ÉAU}
\newcommand{\softrobdescOne}{Rétablissement de la flexion des doigts par actionnement pour des patients en rééducation, évalué par l'amplitude articulaire obtenue, grâce à un exosquelette inspiré de l'origami et actionné par fils en alliage à mémoire de forme (SMA).}
\newcommand{\softrobdescTwo}{Reproduction de la souplesse des tissus humains par association de fils SMA et de polymère Dragon Skin en remplacement de liaisons rigides, avec évaluation du transfert d'effort à travers une interface souple.}

\newcommand{\exoplanettitle}{Simulateur de transits d'exoplanètes}
\newcommand{\exoplanetdates}{Printemps 2025}
\newcommand{\exoplanetlocation}{Abou Dabi, ÉAU}
\newcommand{\exoplanetdescOne}{Simulation de transits d'exoplanètes par implémentation de la mécanique orbitale et de la physique stellaire en Python, avec courbes de lumière conformes à la profondeur et à la durée attendues.}
\newcommand{\exoplanetdescTwo}{Généralisation à différentes configurations étoile--planète grâce à un modèle paramétrable plutôt qu'à des cas codés en dur, avec simulations sur plusieurs jeux de paramètres.}

\newcommand{\windTurbinetitle}{Éolienne flottante en mer -- Stabilisation active}
\newcommand{\windTurbinedates}{05/2026 - 06/2026}
\newcommand{\windTurbinelocation}{Besançon, France}
\newcommand{\windTurbinedescOne}{Garantie de la stabilité passive en tangage d'une plateforme semi-submersible à l'échelle 1:1000, avec hauteur métacentrique GM = 5,38 cm et raideur de rappel K = 0,0607 N$\cdot$m/rad, par calcul hydrostatique d'une géométrie à trois colonnes de type WindFloat sous SolidWorks.}
\newcommand{\windTurbinedescTwo}{Stabilisation des mouvements d'une plateforme semi-submersible en bassin à houle, avec rejet des perturbations en tangage et roulis en boucle fermée sur un prototype instrumenté, grâce à une commande à deux axes avec anti-windup et gestion de la zone morte et de la saturation.}

\newcommand{\carbonmetertitle}{CARBONMETRE -- Écoconception et mesure carbone}
\newcommand{\carbonmeterdates}{09/2025 - 12/2025}
\newcommand{\carbonmeterlocation}{Besançon, France}
\newcommand{\carbonmeterdescOne}{Élaboration d'une méthode de mesure carbone pour une pièce fabriquée, avec procédure traçable des étapes de fabrication à la valeur déclarée, combinant données de fabrication industrielle et pratiques métrologiques.}
\newcommand{\carbonmeterdescTwo}{Mise en relation des choix de fabrication et de leur impact mesuré par comparaison du carbone incorporé de différentes gammes de fabrication, en structurant les données de matériaux et de procédés pour rendre les options directement comparables.}

\newcommand{\hvactitle}{Amélioration thermique d'une chambre en résidence -- Refroidissement passif et évaporatif}
\newcommand{\hvacdates}{06/2026 - 06/2026}
\newcommand{\hvaclocation}{Besançon, France}
\newcommand{\hvacdescOne}{Réduction de la surchauffe d'une chambre en résidence aux heures de pointe, avec une température inférieure de 3 \textdegree{}C à l'extérieur entre 14 h et 21 h, grâce à un écran extérieur composé de film réfléchissant, de PSE de 1 cm et de carton ondulé.}
\newcommand{\hvacdescTwo}{Blocage des apports solaires directs sur une fenêtre à simple vitrage de 9 m$^2$ grâce à des couches réfléchissantes et isolantes extérieures, avec réduction mesurée de la température de surface au pic d'ensoleillement.}

\newcommand{\controlBoxtitle}{Installation électrique complète -- Équipements pour événements de taille moyenne}
\newcommand{\controlBoxdates}{06/2023 - 08/2023}
\newcommand{\controlBoxlocation}{Mexico, Mexique}
\newcommand{\controlBoxdescOne}{Réduction de 30\% du temps de mise sous tension des événements, évalué en heures de montage avant l'ouverture au public, par conception d'un coffret de commande à automate programmable de 12 entrées.}
\newcommand{\controlBoxdescTwo}{Conception d'une installation complète de taille moyenne avec distribution prête pour la mise en service, en dimensionnant les circuits des charges audiovisuelles et de l'alimentation de scène.}

\newcommand{\hyperlooptitle}{NYC Hyperloop -- Véhicule à échelle réduite}
\newcommand{\hyperloopdates}{01/2024 - 05/2024}
\newcommand{\hyperlooplocation}{New York, États-Unis}
\newcommand{\hyperloopdescOne}{Réalisation d'un véhicule Hyperloop à échelle réduite pour un concept de liaison New York--Long Island, avec prototype construit et intégré sous ma responsabilité à la tête du sous-groupe mécanique.}
\newcommand{\hyperloopdescTwo}{Vérification de la lévitation et du comportement aérodynamique par études couplées électromagnétiques et d'écoulement sous COMSOL et ANSYS, dont les résultats ont orienté les évolutions de conception.}

\newcommand{\dronetitle}{Quadricoptère fonctionnel -- Structure allégée}
\newcommand{\dronedates}{03/2024 - 05/2024}
\newcommand{\dronelocation}{New York, États-Unis}
\newcommand{\dronedescOne}{Réduction de 40\% de la masse de la structure par rapport au modèle de référence, mesurée sur les pièces imprimées, grâce à l'impression 3D en PLA et TPU.}
\newcommand{\dronedescTwo}{Respect des critères de vibration et de chute grâce à des itérations géométriques sous SolidWorks et ANSYS, avec résultats d'essai comparés aux seuils de performance fixés.}

\newcommand{\trentxwbtitle}{Rolls-Royce Trent XWB -- Modèle 3D fonctionnel}
\newcommand{\trentxwbdates}{01/2024 - 05/2024}
\newcommand{\trentxwblocation}{New York, États-Unis}
\newcommand{\trentxwbdescOne}{Reproduction du Trent XWB sous forme d'un assemblage fonctionnel complet sous SolidWorks et Fusion 360, couvrant tous les modules par rétro-ingénierie à partir de plans publics.}
\newcommand{\trentxwbdescTwo}{Amélioration de l'efficacité de conception par études d'écoulement et de structure sous ANSYS sur les sections critiques, en évaluant les résultats simulés d'aérodynamique et de déformation.}

\newcommand{\liostitle}{LIOS -- Assistant vocal avec IA}
\newcommand{\liosdates}{12/2024 - Aujourd'hui}
\newcommand{\lioslocation}{Abou Dabi, ÉAU}
\newcommand{\liosdescOne}{Développement d'un assistant vocal opérationnel en Python avec traitement automatique du langage et intégration d'API, permettant l'exécution quotidienne de tâches en mode mains libres.}
\newcommand{\liosdescTwo}{Maintien du contexte sur plusieurs échanges grâce à une mémoire contextuelle conservant les informations de session entre les requêtes.}

\newcommand{\fitpaltitle}{FitPal -- Dispositif intelligent d'entraînement physique}
\newcommand{\fitpaldates}{09/2021 - 12/2021}
\newcommand{\fitpallocation}{Abou Dabi, ÉAU}
\newcommand{\fitpaldescOne}{Codéveloppement d'un dispositif intelligent d'entraînement physique sur M5Stack et de son micrologiciel, avec retour en temps réel sur l'exécution des exercices.}
\newcommand{\fitpaldescTwo}{Transmission en direct des données des capteurs vers une application mobile via Arduino IoT Cloud, grâce à l'intégration des points de connexion IoT.}

\newcommand{\covidtitle}{Cabine de désinfection automatisée}
\newcommand{\coviddates}{09/2020 - 02/2021}
\newcommand{\covidlocation}{Mexico, Mexique}
\newcommand{\coviddescOne}{Réalisation d'un prototype fonctionnel de cabine de désinfection automatisée pilotée par Arduino, intégrant capteurs, brumisation et commande sans fil.}
\newcommand{\coviddescTwo}{Réduction de l'exposition de l'opérateur par activation sans fil et sans contact du cycle de désinfection.}

\newcommand{\rovertitle}{Rover téléopéré -- Numérisation spatiale 3D}
\newcommand{\roverdates}{09/2026 - Aujourd'hui}
\newcommand{\roverlocation}{Gijón, Espagne}
\newcommand{\roverdescOne}{Mise en place d'une acquisition 3D de bout en bout depuis une plateforme télécommandée, avec acquisition de nuages de points sans opérateur dans la zone de numérisation grâce à l'intégration de la commande embarquée. Conception de l'ensemble du châssis et du système mécanique du rover sous SolidWorks, avec analyse structurelle sous ANSYS.}
\newcommand{\roverdescTwo}{Maintien de l'opérateur hors de l'environnement mesuré grâce à une commande entièrement distante du déplacement et de l'acquisition, par développement de la liaison de téléopération et de l'interface de commande.}

\newcommand{\microsensortitle}{Microcapteur piézorésistif de pression de détonation}
\newcommand{\microsensordates}{05/2026 - 06/2026}
\newcommand{\microsensorlocation}{Besançon, France}
\newcommand{\microsensordescOne}{Conception d'un capteur de pression pour des environnements de détonation à 10 GPa, répondant simultanément aux exigences de tenue et d'interface de 50 $\Omega$, par sélection de l'empilement complet de matériaux sous contraintes électriques, mécaniques et économiques couplées.}
\newcommand{\microsensordescTwo}{Conception d'un microcapteur piézorésistif avec automatisation des itérations de géométrie et de masques par code imposant des contraintes de K, $\Delta R/R_0$ et résistivité ; traduction des spécifications en masque lithographique et en gamme de fabrication à l'échelle du wafer. Note de 20/20 en microfabrication.}

% --- EXPÉRIENCE ---
\newcommand{\carbonorg}{Energy Innovation Capital}
\newcommand{\carbontitle}{Analyste de marché -- Captage du carbone}
\newcommand{\carbondates}{10/2023 - 01/2024}
\newcommand{\carbonlocation}{Abou Dabi, ÉAU}
\newcommand{\carbonbullets}{%
  \item Élaboration de recommandations stratégiques d'investissement sur des jeunes entreprises de captage du carbone, avec sélection présentée aux analystes seniors, à partir d'une comparaison homogène des filières CCUS sur l'ensemble de la chaîne de valeur.
  \item Cartographie du marché CCUS donnant une vision structurée des acteurs du captage, du transport et du stockage, par croisement de la littérature scientifique et des annonces de projets commerciaux.
  \item Distinction entre avancées techniques réelles et arguments marketing par comparaison de la maturité technologique des voies de captage, avec examen des données publiées de performance et de coût.
  \item Transformation des résultats d'ingénierie en notes d'aide à la décision et en synthèses de risques destinées aux analystes sur des actifs de technologies climatiques, en traduisant les conclusions techniques en termes de marché.
}

\newcommand{\spacelaborg}{NYU Abu Dhabi, Centre des sciences spatiales}
\newcommand{\spacelabtitle}{Ingénieur et chercheur R\&D -- Plateforme de sciences planétaires}
\newcommand{\spacelabdates}{05/2025 - 09/2025}
\newcommand{\spacelablocation}{Abou Dabi, ÉAU}
\newcommand{\spacelabbullets}{%
  \item Réduction de 40\% du temps de la chaîne de traitement des données, évalué de bout en bout sur des jeux d'images à l'échelle planétaire, grâce à des algorithmes de compression et de reprojection assistés par apprentissage automatique.
  \item Conception et déploiement d'un atlas de Mars fondé sur des images multispectrales réelles de la mission Emirates Mars Mission HOPE, jusqu'à une \href{https://romvio.github.io/mars-viewer/}{application web et mobile} opérationnelle et une version iOS destinée à l'App Store d'Apple aux ÉAU.
  \item Amélioration de la fidélité des cartes entre passages orbitaux superposés, avec homogénéité de la luminosité et continuité des raccords après assemblage, par calibration des recouvrements et reprojection cylindrique vers Robinson.
  \item Ancrage de l'atlas dans des données physiques par collecte de sols analogues dans le désert des ÉAU et caractérisation de la structure et de la morphologie des grains pour reproduire le régolithe martien.
}

\newcommand{\geoorg}{Global Education, NYU Abu Dhabi}
\newcommand{\geotitle}{Chargé de coordination Global Education -- Étudiants en mobilité entrante}
\newcommand{\geodates}{12/2025 - 05/2025}
\newcommand{\geolocation}{Abou Dabi, ÉAU}
\newcommand{\geobullets}{%
  \item Organisation de programmes immersifs universitaires et culturels aux ÉAU pour des étudiants internationaux en mobilité entrante, avec sorties réalisées selon le calendrier semestriel, par coordination complète des transports, des sites et des autorisations.
  \item Réduction des risques lors des déplacements de groupe, avec protocoles de sécurité exécutés sans incident, par encadrement de l'équipe de coordinateurs et répartition des responsabilités sur chaque site.
  \item Rôle d'interlocuteur unique des étudiants en mobilité, avec résolution de problèmes dès le premier contact grâce à des échanges directs avec les services du campus.
  \item Coordination des structures d'accueil et continuité des relais entre services par centralisation des communications dans un canal unique par groupe.
}

\newcommand{\geoOneorg}{Global Education, NYU Abu Dhabi}
\newcommand{\geoOnetitle}{Chargé de coordination Global Education -- Étudiants de première année}
\newcommand{\geoOnedates}{09/2026 - 12/2024}
\newcommand{\geoOnelocation}{Abou Dabi, ÉAU}
\newcommand{\geoOnebullets}{%
  \item Amélioration de la transition en première année grâce à des séances d'accueil structurées en début de semestre, organisées avec l'équipe Global Education.
  \item Organisation et conduite de visites et de permanences de conseil hebdomadaires pour familiariser les étudiants avec le campus.
  \item Renforcement de l'intégration communautaire par programmation d'activités universitaires et culturelles et d'événements inclusifs pour plusieurs promotions.
  \item Accompagnement continu des étudiants tout au long du semestre par des permanences régulières de conseil entre pairs sans rendez-vous.
}

\newcommand{\hcorg}{Centre de santé, NYU Abu Dhabi}
\newcommand{\hctitle}{Emploi étudiant -- Responsable des réseaux sociaux}
\newcommand{\hcdates}{09/2024 - 05/2025}
\newcommand{\hclocation}{Abou Dabi, ÉAU}
\newcommand{\hcbullets}{%
  \item Renforcement de l'utilisation des ressources de bien-être, évalué par les interactions et la portée des campagnes de santé, grâce à des contenus adaptés aux différents publics étudiants.
  \item Publication régulière conforme au calendrier mensuel de contenus, par création en interne des supports graphiques et des textes.
  \item Réduction du délai de réponse aux demandes étudiantes reçues dans la boîte partagée, en assurant le premier contact et l'orientation des dossiers.
  \item Maintien d'un ton homogène sur l'ensemble des canaux conformément aux consignes de communication clinique du centre de santé.
}

\newcommand{\alpesorg}{Alpes Technologies México}
\newcommand{\alpestitle}{Ingénieur projets junior -- Gestion d'entrepôt et systèmes électriques}
\newcommand{\alpesdates}{06/2024 - 09/2024}
\newcommand{\alpeslocation}{Mexico, Mexique}
\newcommand{\alpesbullets}{%
  \item Amélioration du facteur de puissance des clients jusqu'à 99\%, avec réduction de 15\% des coûts énergétiques et de 25\% des arrêts électriques, par spécification et installation de batteries de condensateurs françaises après analyse des harmoniques et des pics de tension de postes haute tension.
  \item Remplacement d'un inventaire entièrement manuel par un système de gestion d'entrepôt (WMS) sur mesure en HTML/JS et Google Sheets, avec gain de 60\% d'efficacité opérationnelle et de fabrication.
  \item Mise en place d'une traçabilité par lot pour le contrôle qualité, avec couverture QR en temps réel de chaque mouvement de stock, grâce aux mises à jour bidirectionnelles et aux droits d'accès par rôle dans un outil unique.
  \item Remise en état de l'installation électrique d'un complexe hôtelier sinistré par un incendie, avec schémas unifilaires et recommandations correctives acceptés par le conseil de direction du client, grâce à des mesures sur site pendant l'urgence.
}

\newcommand{\epigsaorg}{EPIGSA}
\newcommand{\epigsatitle}{Assistant en ingénierie des données -- SIG et hydraulique}
\newcommand{\epigsadates}{05/2024 - 07/2024}
\newcommand{\epigsalocation}{Mexico, Mexique}
\newcommand{\epigsabullets}{%
  \item Consolidation des données d'infrastructures municipales de traitement des eaux, avec couches SIG validées dans PostgreSQL, par intégration et vérification croisée des sources de relevés.
  \item Réduction de 25\% des erreurs de données, mesurées après validation automatisée, grâce à des scripts d'assurance et de contrôle qualité dans la chaîne d'ingestion.
  \item Création d'un processus d'intégration documenté et reproductible par l'équipe, en normalisant les schémas et les conventions de nommage.
  \item Fourniture d'une référence géométrique vérifiée unique aux projets par rapprochement des enregistrements de sources contradictoires, résolution des doublons et correction des incohérences de projection à la source.
}

\newcommand{\solarshiporg}{Solar Ship}
\newcommand{\solarshiptitle}{Stagiaire en stratégie et conception de mobilité aérienne}
\newcommand{\solarshipdates}{05/2023 - 08/2023}
\newcommand{\solarshiplocation}{Toronto, Canada}
\newcommand{\solarshipbullets}{%
  \item Conception d'une logistique aérienne entièrement électrique pour des communautés hors réseau, avec concept de routes et de charges utiles sur des corridors isolés de Zambie et de Tanzanie, par dimensionnement des systèmes d'aéronefs autour de pistes non desservies.
  \item Intégration de la durabilité au concept d'opérations au moyen de critères d'impact social notés dans l'étude comparative des options, en traitant l'accès des communautés comme une contrainte de conception plutôt qu'un simple résultat.
  \item Contribution à un programme d'ingénierie multinational par livraison de dossiers de conception à l'équipe élargie et harmonisation des hypothèses des sous-systèmes entre disciplines.
  \item Justification des infrastructures de liaisons aériennes isolées par comparaison documentée avec les transports terrestres, en évaluant le temps d'accès et le coût par itinéraire.
}

\newcommand{\crousorg}{CROUS Besançon Bourgogne Franche-Comté}
\newcommand{\croustitle}{Emploi étudiant -- Restauration universitaire}
\newcommand{\crousdates}{10/2024 - Aujourd'hui}
\newcommand{\crouslocation}{Besançon, France}
\newcommand{\crousbullets}{%
  \item Maintien de la cadence en cuisine aux heures de pointe, avec continuité du service de midi, en assurant les postes de plonge et de préparation.
  \item Respect des exigences d'hygiène et réussite des contrôles de sécurité alimentaire grâce à l'application des procédures de nettoyage et de désinfection.
  \item Réduction des interruptions de service par réapprovisionnement pendant l'activité et appui au service en salle.
  \item Fiabilité des prises de poste, avec présence régulière pendant l'année universitaire et adaptation à la rotation des affectations.
}

\newcommand{\rallyorg}{Championnat du monde des rallyes -- Guanajuato}
\newcommand{\rallytitle}{Commissaire technique officiel FIA}
\newcommand{\rallydates}{11/2019 - 04/2020}
\newcommand{\rallylocation}{Guanajuato, Mexique}
\newcommand{\rallybullets}{%
  \item Garantie de la conformité technique, sans véhicule non conforme au départ, par marquage officiel et inspection des pneumatiques selon les procédures ISO 9001.
  \item Réalisation des contrôles avant chaque session dans les délais du week-end de course, en coordination avec les commissaires responsables.
  \item Communication claire des décisions de contrôle, acceptées par les pilotes et les équipes, en précisant le règlement enfreint et la correction nécessaire.
  \item Tenue de registres de marquage traçables et auditables par concurrent, avec saisie des données au moment du contrôle.
}

\newcommand{\papaloteorg}{Papalote Museo del Niño}
\newcommand{\papalotetitle}{Chargé d'accueil et de médiation}
\newcommand{\papalotedates}{05/2019 - 07/2019}
\newcommand{\papalotelocation}{Mexico, Mexique}
\newcommand{\papalotebullets}{%
  \item Présentation d'expositions scientifiques et techniques aux enfants par démonstrations pratiques maintenant l'intérêt des visiteurs dans les salles.
  \item Amélioration de l'expérience de visite par propositions et essais avec l'équipe d'exploitation, avec modifications des espaces mises en œuvre après examen par le personnel.
  \item Accueil d'un public varié en adaptant les explications aux différents âges et besoins.
  \item Maintien d'un fonctionnement sûr des dispositifs par surveillance sans incident et application des règles d'accès aux expositions.
}
